\documentclass[aspectratio=169, t, envcountsect, smaller, notheorems]{beamer}
\usetheme[sectionpage=simple, titleformat=allsmallcaps]{metropolis}
\usepackage[ngerman]{babel}
\usepackage{amssymb, mathtools, cancel, pgfplots, nicefrac, stmaryrd}
\usepackage[ruled, linesnumbered, vlined, ngerman, algosection]{algorithm2e}
\usepackage{array, csquotes, hyperref, tasks, multicol}

\MakeOuterQuote{"}

\newcolumntype{L}[1]{>{\raggedleft\arraybackslash}p{#1}}

\usetikzlibrary{automata, positioning, arrows, patterns, fillbetween}

\title[Algorithmen und Datenstrukturen]{Algorithmen und Datenstrukturen}
\subtitle{Sommersemester 2019}
\author{David Luis Wiegandt}
\institute{Wissensmanagement in der Bioinformatik\\Institut für Informatik\\Humboldt-Universität zu Berlin}
\date{}
\titlegraphic{\includegraphics[width=70pt, height=70pt]{images/husiegel_bw.eps}}

\setsansfont[Path = fonts/, UprightFont = *-Regular, BoldFont = *-Medium, ItalicFont = *-Italic]{FiraSans}

\resetcounteronoverlays{algocf}
\resetcounteronoverlays{theorem}

\settasks{label-format=\bfseries, counter-format=(tsk[a]), label-offset=.6em}

\uselanguage{german}
\languagepath{german}
\deftranslation[to=german]{Solution}{L\"osung}
\deftranslation[to=german]{Example}{Beispiel}
\deftranslation[to=german]{Task}{Aufgabe}
\deftranslation[to=german]{Remark}{Bemerkung}
\deftranslation[to=german]{Corollary}{Korollar}

\definecolor{hublue}{RGB}{0, 55, 108}
\definecolor{lighthublue}{RGB}{9, 80, 159}
\definecolor{tealblue}{RGB}{56, 134, 151}
\definecolor{madderlake}{RGB}{204, 41, 54}
\definecolor{ashgrey}{RGB}{174, 195, 176}
\definecolor{mandarin}{RGB}{239, 123, 69}
\definecolor{amaranth}{RGB}{160, 26, 125}

\newcommand<>{\emphlightblue}[1]{\textit#2{\color#2{lighthublue}#1}}
\newcommand<>{\emphteal}[1]{\textit#2{\color#2{teal}#1}}
\newcommand<>{\emphred}[1]{\textit#2{\color#2{madderlake}#1}}
\newcommand<>{\emphorange}[1]{\textit#2{\color#2{mandarin}#1}}
\newcommand<>{\emphpurple}[1]{\textit#2{\color#2{amaranth}#1}}

\newcommand<>{\lightblue}[1]{{\color#2{lighthublue}#1}}
\newcommand<>{\teal}[1]{{\color#2{teal}#1}}
\newcommand<>{\red}[1]{{\color#2{madderlake}#1}}
\newcommand<>{\orange}[1]{{\color#2{mandarin}#1}}
\newcommand<>{\purple}[1]{{\color#2{amaranth}#1}}

\makeatletter
\patchcmd{\beamer@sectionintoc}{\vfill}{\vskip\itemsep}{}{}

\setbeamertemplate{frame footer}{\insertshorttitle}
\setbeamertemplate{itemize item}{\fontsize{3}{4}\raise2.5pt\hbox{\donotcoloroutermaths$\blacksquare$}}
\setbeamertemplate{itemize subitem}{\fontsize{2.5}{3.5}\raise2.5pt\hbox{\donotcoloroutermaths$\blacksquare$}}
\setbeamertemplate{itemize/enumerate body begin}{\normalsize}
\setbeamertemplate{itemize/enumerate subbody begin}{\small}
\setbeamertemplate{section in toc}[sections numbered]
\setbeamertemplate{subsection in toc}[subsections numbered]
\setbeamertemplate{theorems}[numbered]

\setbeamercolor{alerted text}{fg=lighthublue}
\setbeamercolor{block title alerted}{fg=madderlake}
\setbeamercolor{footline}{fg=hublue}
\setbeamercolor{frametitle}{fg=hublue, bg=black!2}
\setbeamercolor{normal text}{fg=hublue}
\setbeamercolor{progress bar}{fg=hublue}
\setbeamercolor{title page header}{fg=white, bg=hublue}

\setbeamerfont{alerted text}{shape=\itshape}
\setbeamerfont{author}{series=\bfseries, size=\small}
\setbeamerfont{footline}{shape=\itshape}
\setbeamerfont{page number in head/foot}{size=\scriptsize}
\setbeamerfont{section title}{size=\LARGE}
\setbeamerfont{title}{size=\LARGE}

\setbeamertemplate{title page}{
\vspace*{-1pt}
\begin{beamercolorbox}[wd=\paperwidth,ht=0.618\paperheight]{title page header}
    \vfill
    \par
    \raggedright%
    \leftskip28pt%
    {\usebeamerfont{title} tutorium \\ \vspace{5pt} \inserttitle}
    \par
    \vspace*{3pt}
    {\usebeamerfont{subtitle}\insertsubtitle}
    \vspace*{-5pt}
    \vfill
    \tikz[remember picture, overlay] \node[anchor=center] at (0.8\paperwidth,0) {\inserttitlegraphic};
\end{beamercolorbox}
\vspace*{2pt}
{\usebeamerfont{author}\beamer@shortauthor}

\medskip
\fontsize{10}{8}\selectfont
\par

{\usebeamerfont{institute}\insertinstitute}
\vfill
}
\makeatother

% default styles for plots
\pgfplotsset{width=\textwidth, compat=1.9}
\pgfplotscreateplotcyclelist{hubcols}{
    {tealblue},
    {madderlake},
    {ashgrey}
}
\pgfplotsset{
    tikzDefaults/.style={
        cycle list name=hubcols,
        font=\small,
        line width = 1pt,
        legend style={
            at={(0.5,-0.2)},
            anchor=north, 
            column sep=7pt,
            draw=none,
            fill=black!2
        },
        legend columns=2,
        axis x line=bottom,
        axis y line=middle,
        tick style={draw=hublue, semithick},
        tick align=center
        every axis/.append style={font=\tiny},
    },
    samples=100
}

\tikzset{
    node distance=2cm,
    every state/.style={
        thick,
        fill=lighthublue!7
    },
    initial text={},
    double distance=2pt,
    every edge/.style={
        draw,
        ->,
        >=stealth,
        auto,
        thick,
        text=hublue
    }
}

% remove rules from algos while retaining the ruled layout
\setlength{\algoheightrule}{0pt}
\setlength{\algotitleheightrule}{0pt}

\let\leq=\leqslant
\let\geq=\geqslant

% general math
\newcommand{\nat}{\mathbb{N}}
\newcommand{\natpos}{\mathbb{N}_{\geq 1}}
\newcommand{\real}{\mathbb{R}}
\newcommand{\realpos}{\mathbb{R}_{\geq 0}}
\newcommand{\powerset}{\mathcal{P}}
\newcommand{\bigO}{\mathcal{O}}
\DeclareMathOperator{\cupdot}{\mathbin{\mathaccent\cdot\cup}}
\DeclareMathOperator{\subsetfin}{\mathbin{\subset_{\mathsf{e}}}}
\newcommand{\defiff}{\mathrel{\vcentcolon\Leftrightarrow}}
\DeclarePairedDelimiter\norm{\lVert}{\rVert}

% query functions
\DeclareMathOperator{\adom}{adom}
\DeclareMathOperator{\vars}{var}
\DeclareMathOperator{\cons}{adom}
\DeclareMathOperator{\free}{frei}
\newcommand{\qfn}[1]{[\![#1]\!]}

% algorithms
\SetKwInOut{Input}{input}
\SetKwInOut{Output}{output}
\ResetInOut{output}
\SetKw{True}{true}
\SetKw{False}{false}

% theorem environments
\theoremstyle{definition}
\newtheorem{theorem}[algocf]{\translate{Theorem}}
\newtheorem{lemma}[algocf]{\translate{Lemma}}
\newtheorem{definition}[algocf]{\translate{Definition}}
\newtheorem{example}[algocf]{\translate{Example}}
\newtheorem{task}[algocf]{\translate{Task}}
\newtheorem*{solution}{\translate{Solution}}
\newtheorem*{remark}{Bemerkung}

% enumerate environments
\newenvironment{subtasks}{
    \setbeamertemplate{enumerate item}{(\alph{enumi})}
    \setbeamerfont{enumerate item}{series=\bfseries}
    \begin{enumerate}
}{
    \end{enumerate}
}

\begin{document}
    \maketitle
    
    \begin{frame}{Organisatorisches}
    \begin{block}{Zeit \& Ort}
        \begin{itemize}
            \item Mittwoch (RUD 25, 4.112) und Donnerstag (RUD 25, 4.113) von 13 bis 15 Uhr
            \item Teilnahme ist \alert{freiwillig} und \alert{flexibel}
            \begin{itemize}
                \item Ihr könnt stets an \alert{einem}, \alert{beiden} oder \alert{keinem} der Termine erscheinen
            \end{itemize}
        \end{itemize}
    \end{block}
    
    \begin{block}{Inhalte}
        \begin{itemize}
            \item Fokus ist die \alert{Wiederholung} von Inhalten aus Vorlesung und Übung
            \item Dar\"uber hinaus ist auch die \alert{Vertiefung} von Inhalten aus der Vorlesung m\"oglich
        \end{itemize}
    \end{block}
    
    \begin{block}{Bei Fragen}
        \begin{itemize}
            \item Im Tutorium oder per Mail an \alert{\href{mailto:wigandtd@hu-berlin.de}{wigandtd@hu-berlin.de}}
        \end{itemize}
    \end{block}
    \end{frame}
    
    \begin{frame}{Behandelte Themen}
        \tableofcontents
    \end{frame}
    
    \section{Laufzeitanalyse}

\begin{frame}{Hintergrund}
\begin{block}{Knappheit von Ressourcen}
    \begin{itemize}
        \item Speicherplatz (Register, RAM, HDD/SSD, Tape, ...)
        \item \alert{Rechenleistung} (Taktfrequenz)
    \end{itemize}
\end{block}
\begin{block}{Eine Problemstellung -- viele Lösungsansätze}
    \begin{itemize}
        \item Welchen Ansatz sollen wir nutzen?
        \item Ziel: Klassifiziere Lösungsansätze nach ihrer \alert{Eignung}
        \begin{itemize}
            \item Eignung $\approx$ Ressourcenbedarf
        \end{itemize}
    \end{itemize}
\end{block}
\begin{block}{Berechnungsmodell}
    \begin{itemize}
        \item Wir analysieren \alert{abstrakte Algorithmen} -- folglich st\"utzen wir unsere Analysen auf \alert{abstrakte Maschinen}
    \end{itemize}
\end{block}
\end{frame}

% \begin{frame}{Berechnungsmodell}
% Wir analysieren \alert{abstrakte Algorithmen} -- folglich st\"utzen wir unsere Analysen auf \alert{abstrakte Maschinen}.

% \medskip

% \begin{columns}[T,onlytextwidth]
% \begin{column}{0.54\textwidth}
% Beispielsweise ist die \alert{Random Access Machine} eine zur \alert{Turingmaschine} \"aquivalente und f\"ur unsere Zwecke leichter zu handhabende Abstraktion einen realen Computers.
% \end{column}
% \begin{column}{0.4\textwidth}
% \begin{tikzpicture}
% \foreach \i in {1,...,4}
% \draw[draw=hublue] (0,3.5-\i*0.5) rectangle (0.8,4-\i*0.5) node[pos=.5] {$I_{\i}$};
% \draw[draw=none, anchor=mid] (0,0.5) rectangle (0.8,1.5) node[pos=.5] {$\vdots$};
% \draw[draw=hublue] (0,0) rectangle (0.8,0.5) node[pos=.5] {$I_m$};

% \foreach \i in {1,...,6}
% \draw[draw=hublue] (1+\i*0.5,3) rectangle (1.5+\i*0.5,3.5) node[pos=.5] {$r_{\i}$};
% \draw[draw=none] (4.5,3) rectangle (5,3.5) node[pos=.5] {$\dots$};
% \end{tikzpicture}
% \end{column}
% \end{columns}
% \end{frame}

\begin{frame}<1-2>[label=first_example]{Ein einfaches Beispiel}
\begin{example}
Gegeben sei eine Liste natürlicher Zahlen $a_1, \dots, a_n \in \nat$.

Maximiere $x = a_i \cdot a_j$ mit $1 \leq i < j \leq n$.

\begin{quote}
    In Worten: Finde das \alert{gr\"o{\ss}te} Produkt, das sich aus zwei \alert{verschiedenen} Elementen der Liste bilden lässt.
\end{quote}

\begin{block}<2->{M\"ogliche Ans\"atze}
    \begin{enumerate}
        \item Betrachte alle möglichen Paare und berechne deren Produkt
        \pause
        \item<4-> Suche die beiden gr\"o{\ss}ten Elemente $a_i, a_j$ und gib $a_i \cdot a_j$ aus
    \end{enumerate}
\end{block}
\end{example}
\end{frame}

\begin{frame}
\frametitle<1>{Algorithmus}
\frametitle<2>{Algorithmus \& Analyse}
\begin{columns}[T,onlytextwidth]
    \begin{column}{0.43\textwidth}
        \vspace{-5pt}
        \begin{algorithm}[H]
        	\caption{Maximales Produkt 1}
        	\label{alg:max_prod}
        	\DontPrintSemicolon
        	\Input{$a_1, \dots, a_n \in \nat$}
            \Output{$\max \{ a_i \cdot a_j : 1 \leq i < j \leq n \}$}
        	$x \gets 0$\;
        	\For{$i \gets 1$ \KwTo $n - 1$}{
        	    \For{$j \gets i + 1$ \KwTo $n$}{
        	        \If{$a_i \cdot a_j > x$}{
        	            $x \gets a_i \cdot a_j$\;
        	        }
                }
            }
            \Return{$x$}
        \end{algorithm}
    \end{column}
    \begin{column}<2>{0.51\textwidth}
        Die \"au{\ss}ere \textbf{for}-Schleife wird $(n-1)$-mal durchlaufen, die innere \textbf{for}-Schleife $n-1, n-2, \dots, 1$-mal
        \begin{itemize}
            \item Anwendung der \alert{gaußschen Summenformel} liefert $\frac{n \, (n-1)}{2}$
        \end{itemize}
        
        Der Aufwand des \textbf{if}-Statements ist von $n$ \alert{unabh\"angig}
        \begin{itemize}
            \item Multiplikationen, Vergleiche und Zuweisungen ben\"otigen lediglich \alert{konstant} viel Zeit
        \end{itemize}
        
        Insgesamt $\leq 4 \cdot \frac{n \, (n-1)}{2} + 2$ Instruktionen
    \end{column}
\end{columns}
\end{frame}

\againframe<3->{first_example}

\begin{frame}
\frametitle<1>{Algorithmus}
\frametitle<2>{Algorithmus \& Analyse}
\begin{columns}[T,onlytextwidth]
    \begin{column}{0.44\textwidth}
        \vspace{-5pt}
        \begin{algorithm}[H]
        	\caption{Maximales Produkt 2}
        	\label{alg:max_prod_ii}
        	\DontPrintSemicolon
        	\Input{$a_1, \dots, a_n \in \nat$}
            \Output{$\max \{ a_i \cdot a_j : 1 \leq i < j \leq n \}$}
        	$k \gets 0, \quad \ell \gets 0$\;
        	\For{$i \gets 1$ \KwTo $n$}{
        	    \If{$a_i > a_k$}{
        	        $\ell \gets k, \quad k \gets i$\;
    	        }
    	        \ElseIf{$a_i > a_{\ell}$}{
    	            $\ell \gets i$\;
    	        }
            }
            \Return{$a_k \cdot a_{\ell}$}
        \end{algorithm}
    \end{column}
    \begin{column}<2>{0.5\textwidth}
        Die \textbf{for}-Schleife wird genau $n$-mal durchlaufen
        
        \medskip
        
        Innerhalb der \textbf{if}-Statements f\"uhren wir lediglich grundlegende Operationen wie Vergleiche und Zuweisungen durch
        
        \medskip
        
        Insgesamt $\leq 3 \, n + 4$ Instruktionen
    \end{column}
\end{columns}
\end{frame}

\begin{frame}<1-2>[label=second_example]{Ein komplexeres Beispiel}
\begin{example}
Wir nennen eine Sequenz $A = a_1, \dots, a_n \in \nat$ \alert{partitionierbar}, falls  es $k, \ell \in \{1, \dots, n\}$ gibt, sodass $\sum_{i=k}^{\ell} a_i = \frac{1}{2} \cdot  \sum_{i=1}^{n} a_i$.

Entscheide f\"ur eine Sequenz $A = a_1, \dots, a_n$ mit $\{ a_1, \dots, a_n \} = \{1, \dots, n\}$, ob diese partitionierbar ist.

\begin{quote}
    In Worten: Entscheide f\"ur eine \alert{Permutation} der nat\"urlichen Zahlen von $1$ bis $n$, ob es eine \alert{zusammenh\"angende Subsequenz} gibt, deren Summe die \alert{H\"alfte} der Summe der Gesamtsequenz ist.
\end{quote}

\begin{block}<2->{M\"ogliche Ans\"atze}
    \begin{enumerate}
        \item Berechne die Summe aller möglichen Subsequenzen \emph{from scratch}
        \pause
        \item<4-> Nutze \alert{dynamische Programmierung}, um fr\"uhere Teilergebnisse wiederzuverwenden
    \end{enumerate}
\end{block}
\end{example}
\end{frame}

\begin{frame}
\frametitle<1->{Algorithmus}
\frametitle<2->{Algorithmus \& Analyse}
\begin{columns}[T,onlytextwidth]
    \begin{column}{0.43\textwidth}
        \vspace{-5pt}
        \begin{algorithm}[H]
        	\caption{Subsequenz 1}
        	\label{alg:subseq_part_1}
        	\DontPrintSemicolon
        	\Input{$A = a_1, \dots, a_n \in \nat$}
            \Output{Ist $A$ partitionierbar?}
        	{\color<2>{madderlake}$s \gets \frac{n \, (n+1)}{4}$\;}
        	\For{$k \gets 1$ \KwTo $n$}{
        	    \For{$\ell \gets k$ \KwTo $n$}{
        	        {\color<3>{madderlake}$s' \gets 0$\;}
        	        \For{$i \gets k$ \KwTo $\ell$}{
        	            {\color<5>{madderlake}$s' \gets s + a_i$\;}
                    }
                    \If{{\color<3>{madderlake}$s' = s$}}{
                        {\color<4>{madderlake}\Return{\True}}
                    }
                }
            }
            {\color<4>{madderlake}\Return{\False}}
        \end{algorithm}
    \end{column}
    \begin{column}<2->{0.51\textwidth}
        \only<2->{
            Die \red<2>{Initialisierung} ben\"otigt 4 Instruktionen
            \begin{itemize}
                \item Eine Addition, eine Multiplikation, eine Division und eine Zuweisung
            \end{itemize}
        }
        
        \only<3->{
            Die \red<3>{Zuweisung} in Zeile $4$ und der \red<3>{Vergleich} in Zeile $6$ werden h\"ochstens $\frac{n \, (n+1)}{2}$-mal ausgef\"uhrt
        }
        
        \medskip
        
        \only<4->{
            Einmalig wird ein Wert \red<4>{zurückgegeben}
        }
        
        \medskip
        
        \only<5->{
            Die \red<5>{innere \textbf{for}-Schleife} wird $\frac{n \, (n+1) \, (n+2)}{6}$-mal durchlaufen (Details: n\"achste Folie!)
            \begin{itemize}
                \item Entsprechend wird auch Zeile $6$ so oft ausgef\"uhrt
            \end{itemize}
        }
        
        \only<6->{
            Insgesamt $\leq \nicefrac{1}{3} \, n \, (n+1) \, (n+5) + 5$ Instruktionen
        }
    \end{column}
\end{columns}
\end{frame}

\begin{frame}{Details zur Laufzeitschranke}
    \begin{block}{Beobachtungen}
        \begin{itemize}
            \item Die innere \textbf{for}-Schleife wird stets $(\ell - k + 1)$-mal durchlaufen
            \begin{itemize}
                \item F\"ur $k = 1$ also $1, \dots, n$-mal
                \item F\"ur $k = 2$ also $1, \dots, (n-1)$-mal
                \item F\"ur $k = n$ also $1, \dots, (n-n+1)$-mal
            \end{itemize}
            \item Als Summe ergibt sich $\sum_{i=1}^n \frac{i \, (i+1)}{2} = \frac{1}{2} \, \sum_{i=1}^n i \, (i+1) = \frac{1}{2} \, \left(\sum_{i=1}^n i^2 + \sum_{i=1}^n i \right)$
            \begin{itemize}
                \item Wir wissen, dass $\sum_{i=1}^n i = \frac{1}{2} \, n \, (n+1)$
                \item F\"ur $\sum_{i=1}^n i^2$ k\"onnen wir per vollst\"andiger Induktion zeigen, dass $\sum_{i=1}^n i^2 = \frac{1}{6} \, n \, (n + 1) \, (2 \, n + 1)$ gilt
            \end{itemize}
        \end{itemize}
    \end{block}
\end{frame}

\begin{frame}{Details zur Laufzeitschranke (2)}
    Insgesamt erhalten wir:
    
    \begin{align*}
        \frac{1}{2} \, \left(\sum_{i=1}^n i^2 + \sum_{i=1}^n i \right) &= \frac{1}{2} \, \left(\frac{1}{6} \, n \, (n + 1) \, (2 \, n + 1) + \frac{1}{2} \, n \, (n+1) \right) \\
        &= \frac{1}{2} \cdot \frac{1}{6} \, \left( n \, (n+1) \, (2\,n+1) + 3 \, n \, (n+1) \right) \\
        &= \frac{1}{2} \cdot \frac{1}{6} \, (2\,n+1+3) \, n \, (n+1) \\
        &= \frac{1}{2} \cdot \frac{1}{6} \, 2\, (n+2) \, n \, (n+1) \\
        &= \frac{1}{6} \, n \, (n+1) \, (n+2)
    \end{align*}
\end{frame}

\againframe<3->{second_example}

\begin{frame}{Dynamische Programmierung}
\begin{block}{Grundidee}
    \begin{itemize}
        \item Nutze \alert{Teilergebnisse}, um Ergebnisse f\"ur gr\"o{\ss}ere Instanzen zu bestimmen
        \item Zum Speichern von Teilergebnissen ben\"otigen wir \alert{zus\"atzlichen Speicherplatz}
    \end{itemize}
\end{block}

\begin{example}
Sei $a_1, \dots, a_n \in \nat$ eine Sequenz nat\"urlicher Zahlen.

Um f\"ur alle $1 \leq k \leq n$ die \alert{kumulierte Summe} $s_k = \sum_{i=1}^{k} a_i$ zu bestimmen, k\"onnen wir
\begin{enumerate}
    \item f\"ur jedes $k$ auf's Neue $s_k = a_1 + \cdots + a_{k}$ berechnen, oder
    \item $s_k$ aus $s_{k-1}$ bestimmen, da $s_k = s_{k-1} + a_k$.
\end{enumerate}
\end{example}
\end{frame}

\begin{frame}
\frametitle<1->{Algorithmus}
\frametitle<2->{Algorithmus \& Analyse}
\begin{columns}[T,onlytextwidth]
    \begin{column}{0.43\textwidth}
        \vspace{-5pt}
        \begin{algorithm}[H]
        	\caption{Subsequenz 2}
        	\label{alg:subseq_part_2}
        	\DontPrintSemicolon
        	\Input{$A = a_1, \dots, a_n \in \nat$}
            \Output{Ist $A$ partitionierbar?}
            {\color<2>{madderlake}
                $s \gets \frac{n \, (n+1)}{4}$\;
            	\red<5>{Sei $(m_{i,j}) \in \nat^{n \times n}$ eine Matrix\;}
        	}
        	\For{$k \gets 1$ \KwTo $n$}{
        	    \For{$\ell \gets k$ \KwTo $n$}{
        	        {\color<3>{madderlake}
            	        $m_{i,j} \gets m_{i, j - 1} + a_{\ell}$\;
            	        \If{$m_{i,j} = s$}{
            	            \red<4>{\Return{\True}}
            	        }
        	        }
                }
            }
            \red<4>{\Return{\False}}
        \end{algorithm}
    \end{column}
    \begin{column}<2->{0.51\textwidth}
        Die \red<2>{Initialisierung} ben\"otigt 4 Instruktionen
        \begin{itemize}
            \item Eine Addition, eine Multiplikation, eine Division und eine Zuweisung
        \end{itemize}
    
        \only<3->{
            Gem\"a{\ss} der gaußschen Summenformel wird der \red<3>{Rumpf} der inneren \textbf{for}-Schleife h\"ochstens $\frac{n \, (n+1)}{2}$-mal ausgef\"uhrt
            
            \begin{itemize}
                \item Dieser beinhaltet $4$ Instruktionen -- eine Subtraktion (in $m_{i, j-1}$), eine Addition, eine Zuweisung und ein Vergleich
            \end{itemize}
        }
        
        \only<4->{
            Einmalig wird ein Wert \red<4>{zur\"uckgegeben}
            
            \medskip
        }
        
        \only<5>{
            Insgesamt $2 \, n \, (n+1) + 5$ Instruktionen
            \begin{itemize}
                \item Und \red<5>{$\approx n^2$ zus\"atzlicher Speicherplatz}
            \end{itemize}
        }
        
    \end{column}
\end{columns}
\end{frame}
    \section{\uppercase{$\bigO$}-Notation}

\begin{frame}{Definition}
\begin{columns}[T,onlytextwidth]
\begin{column}{0.47\textwidth}
    \begin{definition}[$\bigO$-Notation]
    Seien $\teal{f}, \red{g}: \nat \to \realpos$ Funktionen.
    Wir schreiben $\teal{f} \in \bigO(\red{g})$ genau dann, wenn es eine Konstante $\orange{c} \in \real_{> 0}$ gibt, sodass ab einem Index $n_0$ f\"ur alle $n \geq n_0$ gilt, dass $\teal{f}(n) \leq \orange{c} \cdot \red{g}(n)$.
    \end{definition}
    
    \medskip
    
    {\footnotesize Das hei{\ss}t, ab einem gewissen Punkt $n_0$ gibt es keinen \alert{wesentlichen} Unterschied mehr zwischen den Funktionswerten von $\teal{f}$ und $\red{g}$.}
\end{column}
\begin{column}{0.47\textwidth}
    \begin{tikzpicture}
        \begin{axis}[tikzDefaults, xmin=-4, xmax=12, ymin=0, ymax=500]
            \foreach \k in {1, ..., 8}{
                \addplot+[domain=-4:12, draw=mandarin!40, line width=0.5pt, forget plot]{1/2 * \k * x^2};
            }
            \addplot+[domain=-4:12, draw=teal]{2*x^2 - 7*x + 18};
            \addlegendentry{$2 \, x^2 - 7 \, x + 18$}
            \addplot+[domain=-4:12, draw=madderlake]{x^2};
            \addlegendentry{$x^2$}
        \end{axis}
    \end{tikzpicture}
\end{column}
\end{columns}
\end{frame}

\begin{frame}{Beispiel}
\begin{example}
Sei \teal{$f(n) \coloneqq 3 \, n + 5$} und \red{$g(n) \coloneqq 2 \, n$}.

Um zu zeigen, dass $\teal{f} \in \bigO(\red{g})$ ist, w\"ahlen wir \orange{$c=2$} und erhalten
\begin{equation*}
    \teal{\only<1>{3 \, n} \only<2->{\cancel{3 \, n}} + 5} \leq \orange{2} \cdot \red{2 \, n} = \only<1>{3 \, n} \only<2->{\cancel{3 \, n}} + n.
\end{equation*}
\only<2->{F\"ur $n \geq 5$ ist also $\teal{f(n)} \leq \orange{2} \cdot \red{g(n)}$ und somit $n_0 = 5$.}
\end{example}
\end{frame}

\begin{frame}{Bemerkungen}
\begin{remark}<+->
In der Praxis schreibt man h\"aufig $f(n) \in \bigO(g(n))$ anstelle von $f \in \bigO(g)$.

Beispielsweise ist $$\teal{2 \, n^2 + n} \in \bigO(\red{n^2})$$ \"aquivalent zu $$\teal{f} \in \bigO(\red{g}) \text{ f\"ur } \teal{f(n) \coloneqq 2 \, n^2 + n} \text{ und } \red{g(n) \coloneqq n^2}.$$
\end{remark}

\begin{remark}<+->
F\"ur Algorithmus \ref{alg:max_prod} konnten wir zeigen, dass dieser h\"ochstens $4 \cdot \frac{n \, (n-1)}{2} + 2 \in \bigO(n^2)$ Instruktionen ausf\"uhrt.
Demgegen\"uber ben\"otigt Algorithmus \ref{alg:max_prod_ii} nur h\"ochstens $3 \, n \in \bigO(n)$ Instruktionen und ist somit \alert{asymptotisch} effizienter als Algorithmus \ref{alg:max_prod}.
\end{remark}
\end{frame}

\begin{frame}{Eine leichte Aufgabe}
\begin{task}<1->
Sei $f(n) \coloneqq 6 \, n + 7$.

Zeigen Sie, dass $f(n) \in \bigO(n)$ ist, d.h. finden Sie $c \in \real_{> 0}, n_0 \in \natpos$, sodass f\"ur alle $n \geq n_0$ gilt, dass $f(n) \leq c \cdot n$.
\end{task}

\begin{solution}<2->
F\"ur $c = 7$ ist $6 \, n + 7 \leq 7 \, n = 6 \, n + n$ f\"ur alle $n \geq n_0 \coloneqq 7$.
\end{solution}
\end{frame}

\begin{frame}{Mehr Aufgaben}
\begin{task}
Beweisen oder widerlegen Sie:
\begin{subtasks}
    \item $12 \, n - 4 \in \bigO(n)$
    \item $7 \, n + 1 \in \bigO(n^2)$
    \item $\log(n^2) \in \bigO(\log n)$
\end{subtasks}
\end{task}
\end{frame}

\begin{frame}{Alternative Definition}
\begin{definition}[$\bigO$-Notation]
Seien $\teal{f}, \red{g}: \real \to \real$ Funktionen.
Wir schreiben $\teal{f} \in \bigO(\red{g})$ genau dann, wenn $\limsup_{n \to \infty} \frac{\teal{f(n)}}{\red{g(n)}} < \infty$ und $\red{g(n)} > 0$ f\"ur alle $n \in \nat$.
\end{definition}

Das folgende Theorem liefert ein bei dieser Charakterisierung h\"aufig sehr n\"utzliches Werkzeug.

\begin{theorem}[Satz von L'Hospital]
...
\end{theorem}
\end{frame}

\begin{frame}{Ein Beispiel zum Satz von L'Hospital}
\begin{example}

\end{example}
\end{frame}

\begin{frame}{Aufgaben zum Satz von L'Hospital}
\begin{task}

\end{task}
\end{frame}

\begin{frame}{Aufgabe zur Analyse von Algorithmen}
\begin{task}

\end{task}
\end{frame}

    \section{Stacks \& Queues}

\begin{frame}{Abstrakte Datentypen}
\begin{block}{Grundlagen}
\begin{itemize}
    \item Ein \alert{abstrakter Datentyp (ADT)} beschreibt eine Menge von Operationen auf einer Menge von Werten
    \item Wir interessieren uns nur f\"ur die \alert{Signaturen} dieser Operationen
\end{itemize}
\end{block}

\begin{block}{Stacks \& Queues}
\begin{itemize}
    \item Eine \alert{Queue} (Warteschlange) ist ein ADT nach dem \alert{FIFO} (first-in, first-out) Prinzip
    \begin{itemize}
        \item Typische Operationen sind \texttt{enqueue} und \texttt{dequeue}
    \end{itemize}
    \item Ein \alert{Stack} (Stapel) ist ein ADT nach dem \alert{LIFO} (last-in, first-out) Prinzip
    \begin{itemize}
        \item Typische Operationen sind \texttt{push} und \texttt{pop}
    \end{itemize}
    \item Oft steht eine weitere Operation \texttt{head} zur Verf\"ugung, um das "n\"achste" Element zu betrachten, ohne den \alert{Zustand} der Datenstruktur zu modifizieren
\end{itemize}
\end{block}
\end{frame}

\begin{frame}[label=logexpr_intro]{Logische Ausdr\"ucke}
Im Folgenden f\"uhren wir \alert{Syntax} und \alert{Semantik} logischer Ausdr\"ucke ein
\begin{itemize}
    \item Die \alert{Syntax} beschreibt, \alert{wie} ein logischer Ausdruck aufgebaut sein muss
    \item Die \alert{Semantik} beschreibt, \alert{was} ein logischer Ausdruck aussagt
\end{itemize}

\begin{example}<2->
Wie wir gleich sehen werden, ist $\logic{(1 \land (1 \lor 0))}$ ein logischer Ausdruck.
\begin{itemize}
    \item Hierbei nehmen $\logic{1}$ und $\logic{0}$ die Rollen von "\emphteal{wahr}" und "\emphteal{falsch}" ein
    \item Um ihn "auszurechnen", m\"ussen wir u.a. festlegen, was das Ergebnis der Operatoren $\logic{\land}$ und $\logic{\lor}$ sein soll
    \item Intuitiv ist $\logic{(1 \land (1 \lor 0))} = \logic{(1 \land 1)} = \logic{1}$
\end{itemize}
\end{example}
\end{frame}

\begin{frame}{Syntax logischer Ausdr\"ucke}
\begin{definition}\label{sq:def:logexp_syn}
Wir definieren \emphteal{logische Ausdr\"ucke} als Worte \"uber dem Alphabet $$\Sigma = \{ \overbrace{\logic{0}, \, \logic{1}}^{\text{Ziffern}} \} \cup \{ \overbrace{\logic{\land}, \, \logic{\lor}, \, \logic{\oplus}}^{\text{Operatoren}}\} \cup \{ \overbrace{\logic{(}, \, \logic{)}}^{\text{Klammern}} \}.$$

Die Menge der logischen Ausdr\"ucke $\mathcal{L} \subseteq \Sigma^{\star}$ ist rekursiv wie folgt definiert:
\begin{itemize}
    \item Ziffern sind logische Ausdr\"ucke, d.h. $\{ \logic{0}, \logic{1} \} \subseteq \mathcal{L}$
    \item Falls $\teal{w_1}, \teal{w_2} \in \mathcal{L}$, so sind auch $\teal{(w_1 \land w_2)}, \teal{(w_1 \lor w_2)}, \teal{(w_1 \oplus w_2)} \in \mathcal{L}$
\end{itemize}
\end{definition}
\begin{example}<2->
Beispielsweise sind $\logic{1}$, $\logic{(0 \land 1)}$ und $\logic{(1 \land (1 \lor 0))}$ logische Ausdr\"ucke, w\"ahrend $\mathtt{\red{1 \land 1}}$ und $\mathtt{\red{(1 \land 1 \lor 0)}}$ \emph{keine} logischen Ausdr\"ucke sind.
\end{example}
\end{frame}

\begin{frame}{\"Uberpr\"ufung logischer Ausdr\"ucke}
\begin{task}
Sei $w \in \Sigma^{\star}$ mit $\Sigma = \{ \mathtt{0}, \mathtt{1} \} \cup \{ \mathtt{\land}, \mathtt{\lor}, \mathtt{\oplus} \} \cup \{ \mathtt{(}, \mathtt{)} \}$.
Entwerfen Sie einen Algorithmus, der bei Eingabe von $w$ in Zeit $\bigO(n)$ entscheidet, ob $w \in \mathcal{L}$ ist, wobei $n = |w|$ die L\"ange von $w$ in Zeichen sei.
\end{task}
\begin{solution}<2->
Wir arbeiten uns entsprechend der Klammerung von $w$ durch die Eingabe und nutzen den \alert{Call Stack}, um den Kontext vor Auswertung eines Teilausdrucks zu sichern
\begin{itemize}
    \item Dabei gehen wir gem\"a{\ss} der Syntax logischer Ausdr\"ucke vor, verlangen also, dass entweder $w \in \{ \mathtt{0}, \mathtt{1} \}$ oder $w = (w_1 \star w_2)$ mit $w_1, w_2 \in \mathcal{L}$ und $\star \in \{ \land, \lor, \oplus \}$
\end{itemize}
Beachte, dass wir $w$ dabei \emph{nicht} auswerten.
\end{solution}
\end{frame}

\begin{frame}{Algorithmus}
\begin{algorithm}[H]
	\caption{\"Uberpr\"ufung logischer Ausdr\"ucke}
	\label{sq:alg:check}
	\DontPrintSemicolon
    \SetKwFunction{check}{isExpression}
    \SetKwFunction{fsm}{fsm}
    \SetKwProg{Procedure}{procedure}{}{}
    \begin{multicols}{2}
    \Input{$w = w_1w_2\dots w_n \in \Sigma^{\star}$}
    \Output{$w \in \mathcal{L}$?}
    \BlankLine
    \Procedure{\check{$k$}}{
        \lIf{$w_k \in \{ \logic{0}, \logic{1} \}$}{
            \Return{$k$}
        }
        \lElseIf{$w_k = \logic{(}$}{
            \Return{\fsm{$k$}}
        }
        \lElse{\Return{$\bot$}}
    }
    \BlankLine
    \Return{\check{$1$} $= n$}\;
    \Procedure{\fsm{$k$}}{
        $s \gets 0$\;
        \For{$i \gets k + 1$ \KwTo $n$}{
            $\ell \gets$ \check{$i$}\;
            \If{$s \in \{ 0, 2\}$ \And $\ell > 0$}{
                $s \gets s + 1, \quad i \gets \ell$\;
            }
            \ElseIf{$s = 1$ \And $w_i \in \{ \logic{\land}, \logic{\lor}, \logic{\oplus} \}$}{
                $s \gets s + 1$\;
            }
            \ElseIf{$s = 3$ \And $w_i = \logic{)}$}{
                \Return{$i$}\;
            }
            \lElse{\Return{$\bot$}}
        }
    }
    \end{multicols}
\end{algorithm}
\end{frame}

% \begin{frame}{Analyse}
% \end{frame}

\begin{frame}{Alternativer Ansatz I}
\begin{block}{Grundlagen}
Wir k\"onnen die Sprache aller logischen Ausdr\"ucke $\mathcal{L}$ leicht durch Angabe einer \alert{kontextfreien Grammatik} spezifizieren
\begin{itemize}
    \item Dementsprechend k\"onnen wir einen \alert{Kellerautomaten} (PDA) angeben, der f\"ur Worte $w \in \Sigma^{\star}$ entscheidet, ob $w \in \mathcal{L}$
    \begin{itemize}
        \item Zu Details siehe Vorlesung \emph{Einf\"uhrung in die Theoretische Informatik}
    \end{itemize}
\end{itemize}
\end{block}

\begin{block}{Idee}<2->
Wir simulieren den Lauf eines PDA algorithmisch
\begin{itemize}
    \item Um einen PDA zu simulieren, ben\"otigen wir einen \alert{Stack}
    \item Die \alert{\"Uberf\"uhrungsfunktion} implementieren wir durch eine \alert{Lookup-Tabelle}
\end{itemize}

Nun k\"onnen wir jedes Zeichen der Eingabe in Zeit $\bigO(1)$ verarbeiten --- insgesamt $\bigO(n)$
\end{block}
\end{frame}

\begin{frame}[fragile]{Alternativer Ansatz II}
\begin{example}
{
\newcommand{\transition}[3]{\logic{#1}\purple{#2}{\color{hublue},}\,\purple{#3}}
\newcommand{\hash}{\text{\#}}
Folgender PDA $M = (Z, \teal{\Sigma}, \purple{\Gamma}, \delta, q_0, \purple{\text{\#}})$ mit $Z=\{q_0, q_1, q_2, q_3, q_4\}$ und $\purple{\Gamma = \{ \text{\#}, L, R \}}$ erkennt genau die logischen Ausdr\"ucke gem\"a{\ss} Definition \ref{sq:def:logexp_syn}:

\begin{figure}[!h]
\centering
\begin{tikzpicture}[node distance=2.2cm, every loop/.style={looseness=5, out=115, in=65}]
    \node[state, initial] (q0) {$q_0$};
    \node[state, right of=q0] (q1) {$q_1$};
    \node[state, right of=q1] (q2) {$q_2$};
    \node[state, right of=q2] (q3) {$q_3$};
    \node[state, right of=q3] (q4) {$q_4$};
    \draw   (q0) edge[loop above] node[align=center]{$\transition{0}{\hash}{\epsilon}$ \\ $\transition{1}{\hash}{\epsilon}$} (q0)
            (q0) edge[above] node{$\transition{(}{\hash}{L}$} (q1)
            (q1) edge[loop above] node{$\transition{(}{L}{LL}$} (q1)
            (q1) edge[above] node[align=center]{$\logic{0}$ \\ $\logic{1}$} (q2)
            (q2) edge[above] node[align=center]{$\transition{\land}{L}{R}$ \\ $\transition{\lor}{L}{R}$ \\ $\transition{\oplus}{L}{R}$} (q3)
            (q3) edge[bend left, above] node[align=center]{$\logic{0}$ \\ $\logic{1}$} (q4)
            (q3) edge[bend left, below] node{$\transition{(}{R}{RL}$} (q1)
            (q4) edge[bend left, below] node[align=center]{$\transition{\land}{L}{R}$ \\ $\transition{\lor}{L}{R}$ \\ $\transition{\oplus}{L}{R}$} (q3)
            (q4) edge[loop above] node{$\transition{)}{R}{\epsilon}$} (q4);
\end{tikzpicture}
\caption{Ein PDA $M$ mit $L(M) = \mathcal{L}$}
\label{sq:img:pda}
\end{figure}
}
\end{example}
\end{frame}

\againframe<3->{logexpr_intro}

\begin{frame}{Semantik logischer Ausdr\"ucke}
\begin{definition}
Wir definieren den \emph{Wert} $\llbracket \teal{w} \rrbracket$ logischer Ausdr\"ucke $\teal{w} \in \mathcal{L}$ rekursiv wie folgt:
\begin{itemize}
    \item $\eval{\logic{0}} = \logic{0}$ und $\eval{\logic{1}} = \logic{1}$
    \item F\"ur logische Ausdr\"ucke $\teal{w_1}, \teal{w_2}$ ist
    \begin{itemize}
        \item $\eval{\teal{(w_1 \land w_2)}} = \eval{\teal{w_1}} \cdot \eval{\teal{w_2}}$
        \item $\eval{\teal{(w_1 \lor w_2)}} = \begin{cases}\logic{1}, & \text{falls } \eval{\teal{w_1}} + \eval{\teal{w_2}} \geq 1 \\ \logic{0}, & \text{sonst.}\end{cases}$
        \item $\eval{\teal{(w_1 \oplus w_2)}} = \begin{cases}\logic{1}, & \text{falls } \eval{\teal{w_1}} \neq \eval{\teal{w_2}} \\ \logic{0}, & \text{sonst.}\end{cases}$
    \end{itemize}
\end{itemize}
\end{definition}
\begin{example}<2->\label{sq:ex:sem}
Sei $\logic{((0 \lor 1) \land (1 \land 0))}$ gegeben.
Wir erhalten $\eval{
\logic{(}
\underbrace{\logic{(0 \lor 1)}}_{\eval{\logic{(0 \lor 1)}} = \logic{1}}
\logic{\land}
\underbrace{\logic{(1 \land 0)}}_{\eval{\logic{(1 \land 0)}} = \logic{0}}
\logic{)}
} = \logic{0}$.
\end{example}
\end{frame}

\begin{frame}{Auswertung logischer Ausdr\"ucke}
\begin{task}
Entwerfen Sie einen Algorithmus, der einen logischen Ausdruck $\teal{w} \in \mathcal{L}$ auswertet, also $\eval{\teal{w}}$ bestimmt.
Die Laufzeit soll in $\bigO(n)$ liegen, wobei $n = |\teal{w}|$ die L\"ange des logischen Ausdrucks in Zeichen sei.

Nebst einer konstanten Anzahl zus\"atzlicher Variablen primitiven Datentyps d\"urfen Sie eine Queue oder einen Stack benutzen.
\end{task}
\begin{solution}<2->
\begin{quote}
    Idee: Wir werten \teal{w} wie in Beispiel \ref{sq:ex:sem} ``von innen nach au{\ss}en'' aus und merken uns Zwischenergebnisse in einem \alert{Stack}.
\end{quote}

Details folgen!
\end{solution}
\end{frame}
    \section{Sortierverfahren}

\begin{frame}{Grundlagen}
\begin{block}{Listen}
\begin{itemize}
    \item Abstrakter Datentyp zur Repr\"asentation geordneter Sequenzen von Datenwerten
    \item Verschiedene Implementierungen m\"oglich, bspw. als \alert{Array} oder durch \alert{Verkettung}
\end{itemize}
\end{block}

\begin{block}{Vereinbarungen}
Unabh\"angig von der konkreten Implementierung nehmen wir stets an, dass
\begin{itemize}
    \item Listen eine endliche L\"ange $n \in \nat$ haben
    \item Elemente einer Liste durch die Positionen $1, \dots, n$ adressiert werden
    %\item Operationen \texttt{insert}, \texttt{delete} und \texttt{elementAt} zur Verf\"ugung stehen
\end{itemize}
Den Zugriff auf das $i$-te Element einer Liste $L$ notieren wir mit $L_i$.
Die Teilliste vom $i$-ten bis zum $j$-ten Element (mit $i \leq j$) bezeichnen wir mit $L[i, j]$.
\end{block}
\end{frame}

\begin{frame}<1-3>[label=general_sort]
\frametitle<1-3>{Allgemeine Sortierverfahren}
\frametitle<4>{Zur Erinnerung: Allgemeine Sortierverfahren}
\begin{definition}[Allgemeines Sortierverfahren]
Sei $L = \alglist{a_1, \dots, a_n}$ eine Liste \"uber dem durch $\leq_D$ linear geordneten Definitionsbereich $D$.

Ein \alert{\red<2>{allgemeines} Sortierverfahren} berechnet eine Liste $L' = \alglist{a_{\pi(1)}, \dots, a_{\pi(n)}}$ f\"ur eine Permutation $\pi$ von $\{ 1, \dots, n\}$, sodass
\begin{enumerate}
    \item $a_{\pi(i)} \leq_D a_{\pi(i+1)}$ f\"ur alle $1 \leq i < n$ und
    \item<2-> \emphred<2>{nebst $\leq_D$ keine weiteren evtl. auf $D$ definierten Operationen genutzt werden}.
\end{enumerate}
\end{definition}

\begin{remark}<3->
Obgleich wir \emph{oft} Listen \"uber den nat\"urlichen Zahlen $\nat$ betrachten, d\"urfen wir \emph{keinen} Gebrauch von Operationen wie etwa Multiplikation oder Test auf \alert{Kongruenz} machen.
\end{remark}
\pause
\end{frame}

\begin{frame}{\"Uberblick \"uber allgemeine Sortierverfahren}
In der Vorlesung wurden unter anderem folgende \emph{allgemeine} Sortierverfahren vorgestellt:
\begin{itemize}
    \item Bubble Sort
    \item Insertion Sort
    \item Merge Sort
    \item Quick Sort
\end{itemize}
Im Folgenden f\"uhren wir \emph{Schreibtischtests} o.g. Verfahren durch.
\end{frame}

\begin{frame}{Bubble Sort}
\begin{columns}[T,onlytextwidth]
\begin{column}{0.4\textwidth}
\begin{algorithm}[H]
	\caption{Bubble Sort}
	\label{sort:alg:bubblesort}
	\DontPrintSemicolon
	\Input{Liste $L$}
    \Output{Sortierte Liste $L'$}
	\For{$i \gets 1$ \KwTo $n$}{
	    \For{$j \gets 1$ \KwTo $n - i$}{
            \If{$L_j \geq L_{j+1}$}{
                $t \gets L_j$\;
                $L_j \gets L_{j+1}$\;
                $L_{j+1} \gets t$\;
            }
        }
    }
    \Return{$L$}\;
\end{algorithm}
\end{column}
\begin{column}{0.55\textwidth}
\begin{task}<2->
Ordnen Sie folgende Liste mittels des nebenstehenden Algorithmus: $$L = \alglist{17, 5, 4, 63, 46, 58, 4}$$
\end{task}
% \begin{solution}<3->

% \end{solution}
\end{column}
\end{columns}
\end{frame}

\begin{frame}{Insertion Sort}
\begin{columns}[T,onlytextwidth]
\begin{column}{0.4\textwidth}
\begin{algorithm}[H]
	\caption{Insertion Sort}
	\label{sort:alg:insertionsort}
	\DontPrintSemicolon
	\Input{Liste $L$}
    \Output{Sortierte Liste $L'$}
	\For{$i \gets 2$ \KwTo $n$}{
	    $key \gets L_i$\;
	    $j \gets i - 1$\;
	    \While{$L_j > key$ \And $j > 0$}{
            $L_{j+1} \gets L_j$\;
            $j \gets j - 1$\;
        }
    }
    \Return{$L$}\;
\end{algorithm}
\end{column}
\begin{column}{0.55\textwidth}
\begin{task}<2->
Ordnen Sie folgende Liste mittels des nebenstehenden Algorithmus: $$L = \alglist{17, 5, 4, 63, 46, 58, 4}$$
\end{task}
% \begin{solution}<3->

% \end{solution}
\end{column}
\end{columns}
\end{frame}

\begin{frame}{Mergen sortierter Arrays}
Seien $A = \alglist{a_1, \dots, a_m}$ und $B = \alglist{a_{m+1}, \dots, a_{m+n}}$ sortierte Listen der L\"ange $m$ bzw. $n$.
Wir wollen nun eine ebenfalls sortierte Liste $L = \alglist{ a_{\pi(1)}, \dots, a_{\pi(m + n)} }$ bestimmen.
\begin{itemize}
    \item Das geht nat\"urlich, indem wir die Listen $A$ und $B$ zu $L' = \alglist{a_1, \dots, a_m, a_{m+1}, \dots, a_{m+n}}$ konkatenieren und $L'$ dann in $\bigO((m + n) \cdot \log(m + n))$ sortieren
    \item Damit nutzen wir aber nicht, dass $A$ und $B$ bereits sortiert sind
\end{itemize}
Stattdessen durchlaufen wir beide Listen "abwechselnd" von vorn nach hinten und f\"ugen den jeweils kleineren Wert ans Ende eines neuen Arrays ein
\begin{itemize}
    \item Nachteil: Wir ben\"otigen $\bigO(m+n)$ zus\"atzlichen Speicherplatz
\end{itemize}
\end{frame}

\begin{frame}{Mergen sortierter Arrays --- Beispiel}
\begin{example}
Sei $A = \alglist{ 2, 3, 5, 7, 11, 12 }$ und $B = \alglist{ 4, 11, 16, 17 }$.

\medskip

\begin{tabular}[t]{lr}
$A$: &
\begin{tabular}[t]{r|L{0.4cm}|}
    \cline{2-2} 
    \only<1>{$\rightarrow$} & 2 \\ \cline{2-2}
    \only<2>{$\rightarrow$} & 3 \\ \cline{2-2}
    \only<3-4>{$\rightarrow$} & 5 \\ \cline{2-2}
    \only<5>{$\rightarrow$} & 7 \\ \cline{2-2}
    \only<6>{$\rightarrow$} & 11 \\ \cline{2-2}
    \only<7-11>{$\rightarrow$} & 12 \\ \cline{2-2}
\end{tabular}
\end{tabular}
\hfill
\begin{tabular}[t]{lr}
$B$: &
\begin{tabular}[t]{r|L{0.4cm}|}
    \cline{2-2} 
    \only<1-3>{$\rightarrow$} & 4 \\ \cline{2-2}
    \only<4-7>{$\rightarrow$} & 11 \\ \cline{2-2}
    \only<8-9>{$\rightarrow$} & 16 \\ \cline{2-2}
    \only<10-11>{$\rightarrow$} & 17 \\ \cline{2-2}
\end{tabular}
\end{tabular}
\hfill
\begin{tabular}[t]{lr}
$L$: &
\begin{tabular}[t]{r|L{0.4cm}|}
    \cline{2-2} 
    \only<1>{$\rightarrow$} & \only<2->{2} \\ \cline{2-2}
    \only<2>{$\rightarrow$} & \only<3->{3} \\ \cline{2-2}
    \only<3>{$\rightarrow$} & \only<4->{4} \\ \cline{2-2}
    \only<4>{$\rightarrow$} & \only<5->{5} \\ \cline{2-2}
    \only<5>{$\rightarrow$} & \only<6->{7} \\ \cline{2-2}
    \only<6>{$\rightarrow$} & \only<7->{11} \\ \cline{2-2}
    \only<7>{$\rightarrow$} & \only<8->{11} \\ \cline{2-2}
    \only<8>{$\rightarrow$} & \only<9->{12} \\ \cline{2-2}
    \only<9>{$\rightarrow$} & \only<10->{16} \\ \cline{2-2}
    \only<10-11>{$\rightarrow$} & \only<11->{17} \\ \cline{2-2}
\end{tabular}
\end{tabular}
\end{example}
\end{frame}

\begin{frame}{Mergen sortierter Arrays --- Algorithmus}
\begin{algorithm}[H]
	\caption{Merge}
	\label{sort:alg:merge}
	\DontPrintSemicolon
	\Input{Sortierte Listen $A = \alglist{a_1, \dots, a_m}$ und $B = \alglist{a_{m+1}, \dots, a_{m+n}}$}
    \Output{Sortierte Liste $L = \alglist{a_{\pi(1)}, \dots, a_{\pi(m+n)}}$}
    \BlankLine
    \begin{multicols}{2}
    $i \gets 1, \quad j \gets m + 1, \quad k \gets 1$\;
    \While{$i \leq m$ \And $j \leq m + n$}{
        \If{$A_i \leq B_j$}{
            $L_k \gets A_i, \quad i \gets i + 1$
        }
        \Else{
            $L_k \gets B_j, \quad j \gets j + 1$
        }
        $k \gets k + 1$\;
    }
    \If{$i > m$}{
        $L[k, m+n] \gets B[j, m+n]$
    }
    \Else{
        $L[k, m+n] \gets A[i, m]$
    }
    \Return{$L$}\;
    \end{multicols}
\end{algorithm}
\end{frame}

\begin{frame}{Merge Sort}
\begin{columns}[T,onlytextwidth]
\begin{column}{0.42\textwidth}
\begin{algorithm}[H]
	\caption{Merge Sort}
	\label{sort:alg:mergesort}
	\DontPrintSemicolon
	\Input{Liste $L$}
    \Output{Sortierte Liste $L'$}
    \SetKwFunction{MergeSort}{MergeSort}
    \SetKwFunction{Merge}{Merge}
    \SetKwProg{Procedure}{procedure}{}{}
    \Procedure{\MergeSort{$L$}}{
        $n \gets |L|, \quad m \gets \floor{\nicefrac{n}{2}}$\;
        \BlankLine
        \lIf{$n \leq 1$}{
            \Return{$L$}
        }
        \BlankLine
        $L_{\ell} \gets$ \MergeSort{$L[1, m]$}\;
        $L_r \gets$ \MergeSort{$L[m + 1, n]$}\;
        \BlankLine
        \Return{\Merge{$L_{\ell}$, $L_r$}}\;
    }
    \Return{\MergeSort{$L$}}\;
\end{algorithm}
\end{column}
\begin{column}{0.53\textwidth}
\begin{task}<2->
Ordnen Sie folgende Liste mittels des nebenstehenden Algorithmus: $$L = \alglist{17, 5, 4, 63, 46, 58, 4}$$
\end{task}
% \begin{solution}<3->

% \end{solution}
\end{column}
\end{columns}
\end{frame}

\begin{frame}{Partitionierung von Listen}
\begin{definition}[Partitionierung]
Wir nennen eine Liste $L = \alglist{a_1, \dots, a_n}$ \alert{partitioniert bzgl. $a_k$} f\"ur ein $k \in \{ 1, \dots, n \}$, falls $a_i \leq a_k$ f\"ur alle $i < k$ und $a_j \geq a_k$ f\"ur alle $j > k$:
$$L = \alglist{\underbrace{a_1, \dots, a_{k-1}}_{\leq a_k}, a_k, \underbrace{a_{k+1}, \dots, a_n}_{\geq a_k}}$$
\end{definition}

\begin{example}<2->
Die Liste $L = \alglist{4, 7, 3, \teal{8}, 12, 9, \red{8}, 11}$ ist partitioniert bzgl. $\teal{8}$, \emph{nicht} aber bzgl. $\red{8}$.
\end{example}
\end{frame}

\begin{frame}{Partitionierung von Listen --- Problemstellung}
\begin{block}{Problemstellung}
Gegeben eine Liste $L = \alglist{a_1, \dots, a_n}$ und ein $k \in \{ 1, \dots, n \}$.
Bestimme in Zeit $\bigO(n)$ eine Liste $L' = \alglist{a_{\pi(1)}, \dots, a_{\pi(n)}}$, die bzgl. $a_{\pi^{-1}(k)}$ partitioniert ist.
\end{block}

\begin{remark}<2->
Da wir sp\"ater auf den Teillisten $L'[1, \pi^{-1}(k)-1]$ und $L'[\pi^{-1}(k)+1, n]$ links und rechts von Position $\pi^{-1}(k)$ rekursiv absteigen werden, ist idealerweise $\pi^{-1}(k) \approx \nicefrac{n}{2}$.
\end{remark}

\begin{block}{Frage}<3->
Wie k\"onnen wir ein solches $k$ in Zeit h\"ochstens $\bigO(n)$ finden?
\end{block}

\begin{block}{Antwort}<4->
Das ist gleichbedeutend damit, den \alert{Median} von $L$ in Zeit $\bigO(n)$ zu bestimmen --- das ist m\"oglich, aber nicht trivial (siehe bspw. \href{https://www.cs.hhu.de/fileadmin/redaktion/Fakultaeten/Mathematisch-Naturwissenschaftliche_Fakultaet/Informatik/Computational_Cell_Biology/AlDat/Alg2016_03.pdf}{\alert{hier}})

Eine zuver\"assige \alert{Heuristik} ist es, einen ``zuf\"alligen'' Wert zu w\"ahlen --- bspw. immer $a_n$
\end{block}
\end{frame}

\begin{frame}{Partitionierung von Listen --- Algorithmus}
\begin{algorithm}[H]
	\caption{Partitionierung}
	\label{sort:alg:partition}
	\DontPrintSemicolon
	\SetKwFunction{Swap}{swap}
	\SetKwRepeat{Do}{do}{while}
	\Input{Liste $L = \alglist{a_1, \dots, a_n}$}
    \Output{Bzgl. $a_n$ partitionierte Liste $L' = \alglist{a_{\pi(1)}, \dots, a_{\pi(n)}}$}
    $i \gets 1, \quad j \gets n - 1$\;
    \Do{$i < j$}{
        \lWhile{$L_i \leq L_n$ \And $i < n$}{$i \gets i + 1$}
        \lWhile{$L_j \geq L_n$ \And $j > 1$}{$j \gets j - 1$}
        \lIf{$i < j$}{\Swap{$L_i$, $L_j$}}
    }
    \BlankLine
    \Swap{$L_i$, $L_n$}\;
    \BlankLine
    \Return{$L[1, i - 1], L[i + 1, n]$}\;
\end{algorithm}
\end{frame}

\begin{frame}{Quick Sort}
\begin{columns}[T,onlytextwidth]
\begin{column}{0.42\textwidth}
\begin{algorithm}[H]
	\caption{Quick Sort}
	\label{sort:alg:quicksort}
	\DontPrintSemicolon
	\Input{Liste $L$}
    \Output{Sortierte Liste $L'$}
    \SetKwFunction{QuickSort}{QuickSort}
    \SetKwFunction{Merge}{Merge}
    \SetKwProg{Procedure}{procedure}{}{}
    \Procedure{\QuickSort{$L$}}{
        \lIf{$|L| \leq 1$}{\Return{$L$}}
        $L^L, L^R \gets \text{partitioniere } L \text{ bzgl. } L_n$\;
        $L' \gets $ \QuickSort{$L^L$}\;
        $R' \gets $ \QuickSort{$L^R$}\;
    }
    \Return{$\alglist{L', L_n, R'}$}\;
\end{algorithm}
\end{column}
\begin{column}{0.53\textwidth}
\begin{task}<2->
Ordnen Sie folgende Liste mittels des nebenstehenden Algorithmus: $$L = \alglist{17, 5, 4, 63, 46, 58, 4}$$
\end{task}
% \begin{solution}<3->

% \end{solution}
\end{column}
\end{columns}
\end{frame}

% \begin{frame}{Problemreduktion}
% Problemreduktionen sind ein in der theoretischen Informatik h\"aufig genutztes Werkzeug, um per \alert{Widerspruch} zu zeigen, dass ein Problem \emph{nicht} in einer bestimmten Komplexit\"atsklasse liegt.

% \begin{block}{Grundidee}<2->
% Sei $a$ eine Eingabe f\"ur ein Problem $\mathcal{A}$ und analog $b$ eine Eingabe f\"ur ein Problem $\mathcal{B}$.

% Wenn wir 
% \begin{enumerate}
%     \item einen Algorithmus $\mathbb{B}$ haben, der eine Eingabe $b$ f\"ur $\mathcal{B}$ verarbeiten kann (kurz: $\mathbb{B}(b)$ \emph{und} 
%     \item einen Algorithmus $\mathbb{T}$, der eine Eingabe $a$ f\"ur $\mathcal{A}$ in eine Eingabe f\"ur $\mathcal{B}$ \alert{transformieren} kann,
% \end{enumerate}
% dann k\"onnen wir $b \coloneqq \mathbb{F}(a)$ als Eingabe f\"ur $\mathbb{B}$ nutzen und so 

% \end{block}

% \begin{example}<3->
% Siehe Tafel!
% \end{example}
% \end{frame}

\begin{frame}{Bemerkung zur Laufzeit}
In der Vorlesung wurde folgendes Theorem bewiesen:

\begin{theorem}\label{sort:thm:bound}
Sei $L = \alglist{a_1, \dots, a_n}$ eine Liste mit $a_i \neq a_j$ f\"ur $1 \leq i < j \leq n$.
Um $L$ zu sortieren, ben\"otigt jedes \emph{allgemeine Sortierverfahren} im \emph{Worst Case} $\Omega(n \, \log n)$ Vergleiche.
\end{theorem}

\begin{block}{Bedeutung}
Das bedeutet insbesondere, dass wir f\"ur jeden der vorgestellten Algorithmen Instanzen erzeugen k\"onnen, f\"ur deren Sortierung mindestens $n \, \log n$ Vergleiche ben\"otigt werden.

Das gilt aber nur, wenn wir von \emph{beliebigen} Listen als Eingabe ausgehen.
\begin{itemize}
    \item Andernfalls k\"onnen wir gelegentlich durch zus\"atzliche \emph{Annahmen \"uber die Beschaffenheit} der Liste die untere Schranke f\"ur \emph{beliebige} Listen unterbieten.
\end{itemize}
\end{block}
\end{frame}

\begin{frame}{Sortieren spezieller Arrays}
Gelegentlich besitzen wir zus\"atzliches Wissen \"uber die Beschaffenheit einer eingegebenen Liste $L$.
Beispiele daf\"ur sind etwa:
\begin{itemize}
    \item "$L$ ist so vorsortiert, dass alle Elemente der ersten Hälfte von $L$ kleiner sind als alle Elemente der zweiten Hälfte."
    \item "Die Länge von $L$ ist durch 10 teilbar, und jeweils 10 aufeinanderfolgende Elemente (vom Anfang beginnend) sind aufsteigend sortiert."
\end{itemize}
Hierbei stellen wir uns die Frage, inwiefern sich diese Eigenschaften ggf. zur Reduktion der Laufzeit f\"ur's Sortieren nutzen lassen --- etwa auf $\bigO(n)$.

\begin{alertblock}{Achtung}<2->
Die einzelnen Datenwerte d\"urfen wir aber nach wie vor \emphred{nur miteinander vergleichen}!
\end{alertblock}
\end{frame}

% \begin{frame}{Entscheidungsprobleme}
% \begin{definition}[Entscheidungsproblem]
% Sei $U$ eine Menge und sei $A \subseteq U$.
% Das Problem, für ein $x \in U$ zu \emph{entscheiden}, ob auch $x \in A$ ist, wird das \alert{Entscheidungsproblem} für $A$ genannt.

% Ein Algorithmus $\mathbb{A}: U \to \{ 0, 1 \}$ \alert{entscheidet} $A$, falls $\mathbb{A}(x) = 1 \iff x \in A$.

% Das Entscheidungsproblem f\"ur $A$ liegt in $\bigO(g)$ für eine Funktion $g$, falls ein Algorithmus existiert, der $A$ in Zeit $\bigO(g)$ entscheidet.
% \end{definition}

% \begin{example}<2->
% Sei $U = \nat$ die Menge \emph{aller} nat\"urlichen Zahlen und sei $A = \{ 2 \, x : x \in U \}$ die Menge der \emph{geraden} nat\"urlichen Zahlen.
% Die Frage, ob eine beliebige nat\"urliche Zahl $x \in U$ auch gerade ist (also ob $x \in A$), kann algorithmisch leicht in Zeit $\bigO(1)$ entschieden werden, indem wir $\ceil{\nicefrac{x}{2}} \cdot 2 \overset{?}{=} x$ auswerten und ggf. $1$ bzw. $0$ ausgeben.
% \end{example}

% \end{frame}

% \begin{frame}{Problemreduktion}
% Seien $A \subseteq U_1$ und $B \subseteq U_2$ und sei $A$ \emph{nicht} in $\bigO(g)$ entscheidbar.
% Wir wollen zeigen, dass $B$ \emph{ebenfalls nicht} in $\bigO(g)$ entschieden werden kann.

% Zwecks Erzeugung eines \alert{Widerspruchs} nehmen wir an, wir h\"atten einen Algorithmus $\mathbb{B}$, der $B$ in $\bigO(g)$ entscheidet.

% Wenn wir \emph{zeigen} können, dass wir jedes $a \in U_1$ mittels eines Algorithmus $\mathbb{T}$ in Zeit $\bigO(g)$ in ein $\mathbb{T}(a) \in U_2$ \alert{transformieren} können, sodass $a \in A \iff \mathbb{T}(a) \in B$, dann entscheidet folgender Algorithmus $A$ in Zeit $\bigO(g)$:
% \begin{align*}
%     &b \gets \mathbb{T}(a) \\
%     &\mathbf{return} \  \mathbb{B}(b)
% \end{align*}

% Das steht aber im Widerspruch dazu, dass $A$ \emph{nicht} in Zeit $\bigO(g)$ entschieden werden \emph{kann}.
% \end{frame}

\begin{frame}{Aufgaben zum Sortieren spezieller Arrays I}
\begin{task}
Sei $L$ eine Liste der L\"ange $n$ mit $L_i \leq L_{i + 2}$ f\"ur alle $1 \leq i \leq n - 2$.
Gibt es ein \emph{allgemeines Sortierverfahren}, das $L$ in Zeit $\bigO(n)$ sortiert?
\end{task}

\begin{solution}<2->
Betrachte die Teillisten $L' = \alglist{L_1, L_3, L_5, \dots}$ und $L'' = \alglist{L_2, L_4, L_6, \dots}$.
Diese k\"onnen wir in Zeit $\bigO(n)$ aus $L$ bestimmen und per Voraussetzung sind sie bereits sortiert.
Indem wir nun $L'$ und $L''$ in $\bigO(n)$ \emph{mergen}, haben wir $L$ in $\bigO(n)$ sortiert.
\end{solution}
\end{frame}

\begin{frame}{Aufgaben zum Sortieren spezieller Arrays II}
\begin{task}
Sei $L$ eine Liste der L\"ange $n$ mit $L_{\floor{\nicefrac{n}{2}} + 1} = L_{\floor{\nicefrac{n}{2}} + 2} = \dots = L_n$.
Gibt es ein \emph{allgemeines Sortierverfahren}, das $L$ in Zeit $\bigO(n)$ sortiert?
\end{task}

\begin{solution}<2->
Angenommen, es g\"abe einen Algorithmus $\mathbb{A}$, der \emph{solche} Listen in $\bigO(n)$ sortieren kann.
Dann k\"onnen wir in einer \alert{beliebigen} Liste $L$ der L\"ange $n$ einfach einen beliebigen Wert $a$ ans Ende anh\"angen -- und zwar $n$-mal -- und so in Zeit $\bigO(n)$ eine Liste $L'$ der L\"ange $2 \, n$ erhalten, die die Voraussetzungen von Algorithmus $\mathbb{A}$ erf\"ullt.

Wir sortieren nun $L'$ mit $\mathbb{A}$ in Zeit $\bigO(n)$ und l\"oschen anschlie{\ss}end die $n$ hinzugef\"ugten Vorkommen von $a$ wieder aus der sortierten Liste.

Somit haben wir $L$ in Zeit $\bigO(n)$ sortiert --- Widerspruch zu Theorem \ref{sort:thm:bound}!
\end{solution}
\end{frame}

\againframe<4>{general_sort}

\begin{frame}{Spezielle Sortierverfahren}
H\"aufig haben wir \alert{spezielle Kenntnisse} \"uber unsere Eingabedaten
\begin{itemize}
    \item Das Alter von Kunden in Jahren ist in $[0, 122]$
    \item Die K\"orpergr\"o{\ss}e von Patienten in Zentimetern ist in $[6, 272]$
    \item \emph{Hierarchische Strings}
    \begin{itemize}
        \item Raumnummern an der HU wie etwa \texttt{RUD 25, 4'112}
        \item Standorte in Bibliotheken wie etwa \texttt{ST 270 A149}
    \end{itemize}
    \item UNIX Nutzerrechte wie \texttt{666}
\end{itemize}
Wir k\"onnen hier nicht nur vergleichen, sondern auch die \alert{Struktur} der Datenpunkte ausnutzen.
\end{frame}

\begin{frame}{\"Uberblick \"uber spezielle Sortierverfahren}
In der Vorlesung wurden folgende \emph{spezielle} Sortierverfahren vorgestellt:
\begin{itemize}
    \item Radix Exchange Sort
    \item Bucket Sort
\end{itemize}
Im Folgenden f\"uhren wir \emph{Schreibtischtests} o.g. Verfahren durch.

\medskip

\textbf{Algorithmen folgen!}
\end{frame}

    \section{Suchverfahren}

\begin{frame}{Grundlagen}
\begin{block}{Problemstellung}
Gegeben eine \alert{sortierte} Liste $L = \alglist{a_1, \dots, a_n}$ \"uber einem Wertebereich $D$ und ein beliebiges Element $a \in D$.

Wir wollen die Frage beantworten, ob ein $i \in \{ 1, \dots, n \}$ mit $a_i = a$ existiert.
\end{block}
\begin{block}{Ansatz}<2->
Wir sprechen \"uber \alert{sortierte} Listen, k\"onnen also die \alert{Transitivit\"at} der $\leq_D$-Relation nutzen.

Das bedeutet, dass wir nach Vergleich eines Listenelements $a_i \neq a$ nur noch die Teilliste $L[1, i-1]$ bzw. $L[i+1, n]$ betrachten m\"ussen:
$$L = \alglist{\underbrace{a_1, \dots, a_{i-1}}_{\leq a_i}, a_i, \underbrace{a_{i+1}, \dots, a_n}_{\geq a_i}}$$
\end{block}
\end{frame}

\begin{frame}{Suchverfahren}
In der Vorlesung wurden folgende Suchverfahren vorgestellt:
\begin{itemize}
    \item Binary Search
    \item Fibonacci Search
    \item Interpolation Search
\end{itemize}
Im Folgenden führen wir Schreibtischtests o.g. Verfahren durch.
\end{frame}
    %\section{Amortisierte Analyse}

\begin{frame}{Amortisation in der Wirtschaft}
\resizebox{!}{.8\textheight}{%
\begin{tikzpicture}
\begin{axis}[
    tikzDefaults,
    smooth,
	area style,
    xmin=1,
    xmax=19,
    ticks=none,
    axis x line=middle,
    set layers={background, main}
]
    \addplot+[draw=mandarin, fill=mandarin, fill opacity=0.2] coordinates {
        (1, 289)
        (2, 330)
        (3, 361)
        (4, 393)
        (5, 425)
        (6, 462)
        (7, 62)
        (8, 80)
        (9, 107)
        (10, 141)
        (11, 177)
        (12, 226)
        (13, 279)
        (14, 345)
        (15, 422)
        (16, 512)
        (17, 607)
        (18, 718)
    }
    \closedcycle;
    \addplot+[draw=teal, fill=teal, fill opacity=0.2] coordinates {
        (1, 128)
        (2, 129)
        (3, 126)
        (4, 127)
        (5, 122)
        (6, 127)
        (7, 135)
        (8, 143)
        (9, 148)
        (10, 151)
        (11, 165)
        (12, 169)
        (13, 174)
        (14, 182)
        (15, 196)
        (16, 204)
        (17, 219)
        (18, 222)
    }
    \closedcycle;
    \addplot+[draw=madderlake, fill=madderlake, fill opacity=0.2] coordinates {
        (1, -89)
        (2, -88)
        (3, -95)
        (4, -95)
        (5, -90)
        (6, -90)
        (7, -105)
        (8, -105)
        (9, -104)
        (10, -106)
        (11, -108)
        (12, -109)
        (13, -102)
        (14, -100)
        (15, -102)
        (16, -106)
        (17, -106)
        (18, -105)
    }
    \closedcycle;
\end{axis}
\end{tikzpicture}
}
\end{frame}

\begin{frame}{Selbstorganisierende Listen}

\end{frame}
\end{document}
